% sample.tex — Full demonstration of sudodoc.sty features
% Compile: lualatex sample.tex (run TWICE for TOC + overlays)
\documentclass[12pt]{article}

\usepackage{sudodoc}

\sudodocsetup{
  title={A Comprehensive Study of\\Distributed Systems},
  subtitle={Architecture, Consensus, and Fault Tolerance\\in Modern Infrastructure},
  author={Max Mustermann},
  institution={Technische Universit\"at Berlin\\Fakult\"at f\"ur Elektrotechnik und Informatik},
  date={May 2026},
}

\begin{document}

% ── Title Page (dark theme) ──
\maketitlepage

% ── Table of Contents ──
\maketoc

% ══════════════════════════════════════════════════════════
% CHAPTER 1
% ══════════════════════════════════════════════════════════
\chapterpage{1}{Foundations of Distributed Systems}

\section{Introduction}
Distributed systems form the backbone of modern computing infrastructure.
From cloud services to social networks, the ability to coordinate multiple
autonomous machines is fundamental to scalability and reliability.

\subsection{Motivation and Goals}
The primary motivation for studying distributed systems lies in their
ubiquity. Every time you use a web application, stream a video, or check
your email, you are interacting with a distributed system.

\subsubsection{Why Distributed?}
Centralized systems, while simpler to reason about, face inherent
limitations in scale, fault tolerance, and geographic distribution.
A single point of failure can bring down an entire service.

\subsection{Key Concepts}
The following concepts are essential for understanding distributed systems:

\begin{description}
  \item[Consensus] Agreement among distributed nodes on a single
    data value or system state, despite failures and network delays.
  \item[Quorum] The minimum number of nodes required to agree before
    an operation is considered successful, typically a majority.
  \item[Partition Tolerance] The ability of a system to continue
    operating despite network partitions that isolate subsets of nodes.
\end{description}

% ══════════════════════════════════════════════════════════
% CHAPTER 2
% ══════════════════════════════════════════════════════════
\chapterpage{2}{Consensus Algorithms}

\section{Algorithm Comparison}
Various consensus algorithms offer different trade-offs between
safety guarantees, fault tolerance, and implementation complexity.

\vspace{4pt}
\begin{table}[h]
\centering
\caption{Comparison of Consensus Algorithms}
\begin{tabular}{@{} l c c c @{}}
\toprule
\rowcolorhead
\textcolor{tableHeadText}{\textbf{Algorithm}} &
\textcolor{tableHeadText}{\textbf{Consensus}} &
\textcolor{tableHeadText}{\textbf{Tolerance}} &
\textcolor{tableHeadText}{\textbf{Complexity}} \\
\midrule
Paxos        & Safety    & $f < n/2$       & Moderate \\
\rowcoloralt
Raft         & Safety    & $f < n/2$       & Low      \\
Viewstamped  & Safety    & $f < n/2$       & Moderate \\
\rowcoloralt
PBFT         & Safety    & $f < n/3$       & High     \\
Zab          & Total     & $f < n/2$       & Moderate \\
\rowcoloralt
Tendermint   & Total     & $f < n/3$       & Moderate \\
HotStuff     & Total     & $f < n/3$       & Low      \\
\bottomrule
\end{tabular}
\end{table}

\section{Raft Protocol}
Raft was designed as an understandable alternative to Paxos, decomposing
the consensus problem into independent subproblems: leader election,
log replication, and safety.

\subsection{Leader Election Process}
The leader election follows a clear sequence of steps:

\begin{enumerate}
  \item Follower does not receive heartbeat within election timeout
  \item Follower increments its term and becomes a candidate
    \begin{enumerate}
      \item Candidate votes for itself
      \item Candidate sends RequestVote RPCs to all peers
    \end{enumerate}
  \item Candidate waits for votes from majority of nodes
  \item If majority grants votes, candidate becomes leader
  \item Leader sends periodic heartbeats to maintain authority
\end{enumerate}

\subsection{Log Replication}
Once a leader is established, clients send commands to the leader, which
appends them to its log and replicates them across followers using
AppendEntries RPCs.

\begin{itemize}
  \item Leader appends command to its local log
    \begin{itemize}
      \item Each entry stores command and term number
      \item Entry index monotonically increases
    \end{itemize}
  \item Leader sends AppendEntries to all followers
  \item Followers append entry if log is consistent
  \item Leader commits once majority acknowledges
\end{itemize}

% ══════════════════════════════════════════════════════════
% CHAPTER 3
% ══════════════════════════════════════════════════════════
\chapterpage{3}{Implementation Details}

\section{Code Example}
A minimal Raft-style consensus node implementation in Python:

\begin{lstlisting}[language=Python, caption={Raft node skeleton}]
class RaftNode:
    def __init__(self, node_id: int, peers: list):
        self.id = node_id
        self.peers = peers
        self.state = "follower"
        self.term = 0
        self.log = []
        self.voted_for = None

    def request_vote(self, candidate_term: int):
        """Handle incoming RequestVote RPC."""
        if candidate_term > self.term:
            self.term = candidate_term
            self.voted_for = None
            self.state = "follower"
        return self._grant_vote(candidate_term)
\end{lstlisting}

\section{Key Insight}
\begin{blockquote}
A distributed system is one in which the failure of a computer you
didn't even know existed can render your own computer unusable.
--- Leslie Lamport
\end{blockquote}

\vspace{6pt}

This fundamental observation underscores the importance of fault
tolerance as a first-class design concern in distributed systems.

\begin{lstlisting}[language={}, caption={Minimal configuration}]
[cluster]
node_id = "node-1"
bind_addr = "0.0.0.0:7946"
election_timeout = "1.5s"
heartbeat_timeout = "500ms"
\end{lstlisting}

\end{document}
